% !TEX root = ../main.tex

\section{Pedal Networks}
\label{sec:network-model}

The left of \cref{fig:principle} shows four ways to wire a potentiometer to a stereo plug. In the two common wirings at the top, the wiper sits on the ring or on the tip, and the sleeve is the grounded track end \cite{mission2014expression, nonlinearlabs2026pedals}; less often, the wiper sits on the sleeve. Two-wire pedals connect only tip and sleeve, with the wiper tied to one track end. A mono plug additionally shorts the ring of the jack to its sleeve. For a wiper on the tip, this puts the two track halves in parallel, so a track of resistance $R$ at relative travel $x$ shows $R\,x(1-x)$, which rises and falls again. Potentiometer pedals therefore need a stereo cable.

Any passive resistive network between three terminals is externally equivalent to three resistors meeting in an inaccessible star point (\cref{fig:principle}, center) \cite{kennelly1899equivalence}, and only this equivalent can be recovered from the terminals \cite{borcea2012resistor}. The resistance from one contact to the star point is called an \emph{arm}: $r_\mathrm{T}$, $r_\mathrm{R}$ and $r_\mathrm{S}$ for tip, ring and sleeve. A potentiometer has exactly one arm near \qty{0}{\ohm}, the wiper, and the ratio of the other two arms is its position; a two-wire pedal has one isolated arm. The solver reports the relative arm resistances $q_i = r_i/R_\mathrm{tot}$ with respect to the total resistance $R_\mathrm{tot} = r_\mathrm{T} + r_\mathrm{R} + r_\mathrm{S}$.

\begin{figure}[tb]
  \centering
  % !TEX root = ../../main.tex
% Measurement principle (sec:network-model, sec:pair-measurement): four pedals, their star
% equivalent and the pull matrix of one jack, with the paths of one pair measurement highlighted
% (tip high, ring low, sleeve floating and read). Spans \linewidth; every position follows from the
% grid below.
%
% Columns: the pedals (potentiometer, two jogs for its wires, plug), the star (sleeve end, star
% point, ring end, two jogs for the ring and sleeve wires), the pull matrix (contacts, one column
% per contact, the pull resistors, the labels). A backdrop behind the pedals narrows to the one
% behind the star, a dotted link joins the star's wires and the contacts; both gaps share the
% width that is left. Rows: the pull matrix's seven rows at a fixed pitch, from the high rail
% through the contacts to the low rail. The star's tip end lies on the tip row and its other ends on
% the sleeve row, so its arms are 4/3 of a row pitch long; the sleeve wire runs below, half a row
% pitch under the sleeve row. The pedals sit on every second row.

\begin{tikzpicture}[diagram, diagram title font=\scriptsize,
    diagram name font=\fontsize{7}{8.4}\selectfont, diagram text font=\tiny]
  % ---------------------------------------------------------------- grid
  \pgfmathsetlengthmacro\rowPitch{7mm}
  \pgfmathsetlengthmacro\backdropMargin{2mm}  % around the pedals and the star
  \pgfmathsetlengthmacro\titleSpace{3mm}
  % A pedal: the potentiometer's wires run in lanes into the plug's handle.
  \pgfmathsetlengthmacro\laneGap{2.5mm}
  \pgfmathsetlengthmacro\leadLength{1mm}      % beyond the potentiometer's track ends
  \pgfmathsetlengthmacro\wiperLength{2.5mm}
  \pgfmathsetlengthmacro\jogGap{2mm}
  \pgfmathsetlengthmacro\handleGap{3mm}       % between the last jog and the handle
  % The plug, from the handle to the tip.
  \pgfmathsetlengthmacro\handleWidth{4.5mm}
  \pgfmathsetlengthmacro\handleHeight{2*\laneGap + 2.4mm}
  \pgfmathsetlengthmacro\plugDiameter{4.2mm}
  \pgfmathsetlengthmacro\insulatorDiameter{0.8*\plugDiameter}
  \pgfmathsetlengthmacro\sleeveEnd{\handleWidth + 10mm}
  \pgfmathsetlengthmacro\ringStart{\sleeveEnd + 0.6mm}
  \pgfmathsetlengthmacro\ringEnd{\ringStart + 2.5mm}
  \pgfmathsetlengthmacro\tipStart{\ringEnd + 0.6mm}
  \pgfmathsetlengthmacro\plugLength{\tipStart + 4.5mm}
  % The star: three arms at 120 degrees, the lower two 30 degrees below the horizontal.
  \pgfmathsetlengthmacro\armLength{4/3*\rowPitch}
  \pgfmathsetlengthmacro\armReach{cos(30)*\armLength}
  \pgfmathsetlengthmacro\starWidth{2*\backdropMargin + 2*\armReach + 2*\jogGap}
  % The pull matrix.
  \pgfmathsetlengthmacro\columnPitch{10mm}
  \pgfmathsetlengthmacro\labelWidth{8mm}      % the rail tags and labels right of the matrix
  \pgfmathsetlengthmacro\pedalWidth{\wiperLength + 2*\jogGap + \handleGap + \plugLength}
  % The narrowing backdrop and the link share what is left.
  \pgfmathsetlengthmacro\freeWidth{\linewidth - 2*\backdropMargin - \pedalWidth - \starWidth
    - 5*\columnPitch - \labelWidth}
  \pgfmathsetlengthmacro\funnelLength{0.55*\freeWidth}
  \pgfmathsetlengthmacro\linkLength{0.45*\freeWidth}

  % Columns, left to right.
  \coordinate (pedals west) at (0, 0);
  \coordinate (pot) at (\backdropMargin, 0);
  \coordinate (wiper jog) at ($(pot) + (\wiperLength + \jogGap, 0)$);
  \coordinate (track jog) at ($(wiper jog) + (\jogGap, 0)$);
  \coordinate (plug) at ($(track jog) + (\handleGap, 0)$);
  \coordinate (pedal east) at ($(plug) + (\plugLength, 0)$);
  \coordinate (pedals east) at ($(pedal east) + (\backdropMargin, 0)$);
  \coordinate (star west) at ($(pedals east) + (\funnelLength, 0)$);
  \coordinate (star) at ($(star west) + (\backdropMargin + \armReach, 0)$);
  \coordinate (ring jog) at ($(star) + (\armReach + \jogGap, 0)$);
  \coordinate (sleeve jog) at ($(ring jog) + (\jogGap, 0)$);
  \coordinate (star east) at ($(sleeve jog) + (\backdropMargin, 0)$);
  \coordinate (contacts) at ($(star east) + (\linkLength, 0)$);
  \foreach \contact [count=\column] in {tip, ring, sleeve}
    \coordinate (\contact\space column) at ($(contacts) + (\column*\columnPitch, 0)$);
  \coordinate (pull) at ($(contacts) + (4*\columnPitch, 0)$);
  \coordinate (labels) at ($(contacts) + (5*\columnPitch, 0)$);

  % Rows, top to bottom.
  \foreach \row [count=\index from 0] in {high rail, high switches, tip, ring, sleeve,
      low switches, low rail}
    \coordinate (\row) at (0, {-\index*\rowPitch});
  \foreach \pedal in {0, ..., 3}
    \coordinate (pedal \pedal) at (0, {-2*\pedal*\rowPitch});
  \coordinate (sleeve lane) at ($(sleeve) - (0, \rowPitch/2)$);
  % Both backdrops are centered on the ring row.
  \pgfmathsetlengthmacro\pedalsHalfHeight{3*\rowPitch + \handleHeight/2 + 2.5mm}
  \pgfmathsetlengthmacro\starHalfHeight{1.5*\rowPitch + \backdropMargin}
  \coordinate (pedals top) at ($(ring) + (0, \pedalsHalfHeight)$);
  \coordinate (pedals bottom) at ($(ring) - (0, \pedalsHalfHeight)$);
  \coordinate (star top) at ($(ring) + (0, \starHalfHeight)$);
  \coordinate (star bottom) at ($(ring) - (0, \starHalfHeight)$);
  \coordinate (title row) at ($(pedals top) + (0, \titleSpace)$);

  \tikzset{
    plug body/.style={draw=wire, fill=white, line width=0.6pt},
    contact letter/.style={text=muted, inner sep=0pt, font=\fontsize{4.5}{5}\selectfont},
    % A stereo plug from the west side of its handle, which names the lanes of its contacts.
    pics/plug/.style={code={
      \draw[plug body] (0, -\handleHeight/2) rectangle (\handleWidth, \handleHeight/2);
      \draw[plug body] (\handleWidth, -\plugDiameter/2) rectangle (\sleeveEnd, \plugDiameter/2);
      \fill[wire] (\sleeveEnd, -\insulatorDiameter/2) rectangle (\ringStart, \insulatorDiameter/2);
      \draw[plug body] (\ringStart, -\plugDiameter/2) rectangle (\ringEnd, \plugDiameter/2);
      \fill[wire] (\ringEnd, -\insulatorDiameter/2) rectangle (\tipStart, \insulatorDiameter/2);
      \draw[plug body] (\tipStart, -\plugDiameter/2) -- ({(\tipStart + \plugLength)/2}, -\plugDiameter/2)
        -- (\plugLength, -\plugDiameter/6) -- (\plugLength, \plugDiameter/6)
        -- ({(\tipStart + \plugLength)/2}, \plugDiameter/2) -- (\tipStart, \plugDiameter/2) -- cycle;
      \foreach \contact/\lane in {T/1, R/0, S/-1}
        \node[contact letter] at (\handleWidth/2, \lane*\laneGap) {\contact};
    }},
  }

  % ---------------------------------------------------------------- pedals
  % Each wire jogs into its lane: the wiper's first, the track ends' further right, so that a
  % wire crossing to another lane crosses the others only once.
  \foreach \info/\endOne/\wiper/\endThree [count=\pedal from 0] in {wiper on ring/T/R/S,
      wiper on tip/R/T/S, wiper on sleeve/T/S/R, two-wire/T/S/S} {
    \foreach \contact/\lane in {T/1, R/0, S/-1}
      \coordinate (lane \pedal\space\contact) at ($(plug |- pedal \pedal) + (0, \lane*\laneGap)$);
    \node[resistor vertical] (pot \pedal) at (pot |- pedal \pedal) {};
    \coordinate (wiper \pedal) at ($(pot \pedal.east) + (\wiperLength, 0)$);
    \draw[bus] (pot \pedal.north) -- ++(0, \leadLength) to[jog via=track jog] (lane \pedal\space\endOne);
    \draw[bus] (pot \pedal.south) -- ++(0, -\leadLength) to[jog via=track jog] (lane \pedal\space\endThree);
    \draw[bus] (wiper \pedal) to[jog via=wiper jog] (lane \pedal\space\wiper);
    \draw[flow] (wiper \pedal) -- (pot \pedal.east);
    \ifx\wiper\endThree
      \node[junction] at (track jog |- lane \pedal\space\wiper) {};
    \fi
    \pic at (plug |- pedal \pedal) {plug};
    \node[info, anchor=south west]
      at ($(plug |- pedal \pedal) + (\handleWidth + 0.5mm, \plugDiameter/2 + 0.4mm)$) {\info};
  }

  % ---------------------------------------------------------------- star
  % Each arm leaves the star point at its angle; its label sits outside the star, square to it.
  \coordinate (star point) at ($(star |- sleeve) + (0, \armLength/2)$);
  \foreach \contact/\letter/\angle/\side in {tip/T/90/0, ring/R/-30/60, sleeve/S/210/120} {
    \coordinate (\contact\space end) at ($(star point) + (\angle:\armLength)$);
    \coordinate (\contact\space arm) at ($(star point)!0.5!(\contact\space end)$);
    \draw[bus] (star point) -- (\contact\space end);
    \node[resistor, rotate=\angle] at (\contact\space arm) {};
    \pgfmathtruncatemacro\labelAnchor{\side + 180}
    \node[edge label, anchor=\labelAnchor] at ($(\contact\space arm) + (\side:1.2mm)$)
      {$r_\mathrm{\letter}$};
  }
  \node[junction] at (star point) {};
  % To the contact rows: the ring wire jogs up, the sleeve wire runs below the ring's.
  \draw[bus] (tip end) -- (star east |- tip);
  \draw[bus] (ring end) -- (ring jog |- ring end) |- (star east |- ring);
  \draw[bus] (sleeve end) |- (sleeve jog |- sleeve lane) |- (star east |- sleeve);

  % ---------------------------------------------------------------- pull matrix
  % Each contact's column runs from the high rail through its two switches to the low rail.
  \foreach \contact in {tip, ring, sleeve} {
    \draw[bus] (\contact\space column |- high rail)
      -- ($(\contact\space column |- high switches) + (0, \diagramSwitchReach)$);
    \draw[bus] ($(\contact\space column |- high switches) - (0, \diagramSwitchReach)$)
      -- ($(\contact\space column |- low switches) + (0, \diagramSwitchReach)$);
    \draw[bus] ($(\contact\space column |- low switches) - (0, \diagramSwitchReach)$)
      -- (\contact\space column |- low rail);
  }
  % The drives of the pair measurement: tip high, ring low, sleeve floating.
  \pic at (tip column |- high switches) {switch closed};
  \pic at (ring column |- high switches) {switch open};
  \pic at (sleeve column |- high switches) {switch open};
  \pic[yscale=-1] at (tip column |- low switches) {switch open};
  \pic[yscale=-1] at (ring column |- low switches) {switch closed};
  \pic[yscale=-1] at (sleeve column |- low switches) {switch open};
  \foreach \row in {high switches, low switches}
    \draw[actuator] (tip column |- \row) -- (labels |- \row) node[edge label, anchor=west] {drive};

  % The rails, each through its pull resistor.
  \foreach \rail/\hue/\name/\side in {high rail/railHigh/2.5\,V/above, low rail/railLow/GND/below} {
    \node[rail tag=\hue, anchor=west] (\rail\space tag) at (labels |- \rail) {\name};
    \draw[bus] (tip column |- \rail) -- (\rail\space tag);
    \node[resistor, label={[edge label]\side:$R_\mathrm{p}$}] at (pull |- \rail) {};
    \foreach \contact in {ring, sleeve}
      \node[junction] at (\contact\space column |- \rail) {};
  }

  % The contacts, linked to the arms by the plug and read by the ADC.
  \foreach \contact/\letter in {tip/T, ring/R, sleeve/S} {
    \draw[link] (star east |- \contact) -- (contacts |- \contact);
    \draw[bus] (contacts |- \contact) -- (labels |- \contact) node[edge label, anchor=west] {ADC};
    \node[junction] at (\contact\space column |- \contact) {};
    \node[terminal, label={[edge label]above:\letter}] at (contacts |- \contact) {};
  }

  % ---------------------------------------------------------------- background
  \begin{scope}[on background layer]
    % The pedals' backdrop narrows to the star's.
    \fill[backdrop] (pedals west |- pedals top) -- (pedals east |- pedals top)
      -- (star west |- star top) -- (star east |- star top) -- (star east |- star bottom)
      -- (star west |- star bottom) -- (pedals east |- pedals bottom)
      -- (pedals west |- pedals bottom) -- cycle;
    % The pair measurement: the loop current from the high rail to the low rail through the tip
    % and ring arms, and the sleeve, which reads the star point.
    \draw[highlight=railHigh] (high rail tag.west) -- (tip column |- high rail) -- (tip column |- tip)
      -- (tip end) -- (star point);
    \draw[highlight=railLow] (star point) -- (ring end) -- (ring jog |- ring end) |- (ring column |- ring)
      -- (ring column |- low rail) -- (low rail tag.west);
    \draw[highlight=other] (star point) -- (sleeve end) |- (sleeve jog |- sleeve lane)
      |- (labels |- sleeve);
  \end{scope}

  % ---------------------------------------------------------------- titles
  \coordinate (pedal center) at ($(pot)!0.5!(pedal east)$);
  \coordinate (star center) at ($(star west)!0.5!(star east)$);
  \coordinate (matrix center) at ($(contacts)!0.5!(labels) + (\labelWidth/2, 0)$);
  \foreach \section/\title in {pedal/Pedal Wirings, star/Star Equivalent, matrix/Pull Matrix}
    \node[group title=ink, anchor=base] at (\section\space center |- title row) {\title};
\end{tikzpicture}

  \titledcaption{Measurement Principle}{Left, four pedals wire a potentiometer to the tip (T), ring (R) and sleeve (S) of a stereo plug. Center, each of them is equivalent to the arms $r_\mathrm{T}$, $r_\mathrm{R}$ and $r_\mathrm{S}$ meeting in the star point. Right, the pull matrix of one jack: per contact, one switch of the upper row (drive) connects it through the pull resistor $R_\mathrm{p}$ to the \qty{2.5}{\volt} rail, one of the lower row to ground, and the ADC reads it; the dotted links close when the plug is inserted. Highlighted, a pair measurement with the tip high, the ring low and the sleeve floating, which reads the star point.}
  \label{fig:principle}
\end{figure}
